\Needspace{8\baselineskip}
\Article{191}

\textsc{Theorem.} \textit{If a reduced form $(a,b,-a')$ with determinant $D$ goes by the substitution $\alpha,\beta,\gamma,\delta$ into the reduced form $(A,B,-A')$ with the same determinant, then,} first, \textit{$\frac{\pm\surd D-b}{a}$ lies between $\frac{\alpha}{\gamma}$ and $\frac{\beta}{\delta}$ (provided neither $\gamma$ nor $\delta=0$, that is, both bounds are finite), taking the upper sign when neither bound has the sign opposite to that of $a$ (or, more clearly, when either both have the coefficient's sign, or one has that sign and the other is $=0$), and the lower sign when neither has the same sign as $a$;} second, \textit{$\frac{\pm\surd D+b}{a'}$ lies between $\frac{\gamma}{\alpha}$ and $\frac{\delta}{\beta}$ (provided neither $\alpha$ nor $\beta=0$), taking the upper sign when neither bound has the sign opposite to that of $a'$ (or $a$), and the lower sign when neither has the same sign as $a'$}\HistoricalNote{*)}{Other cases clearly cannot occur, since by the preceding article, because $\alpha\delta-\beta\gamma=\pm1$, the two bounds can neither have opposite signs nor both be $=0$.}.

\textit{Proof.} We have the equations
\[
\begin{array}{rcll}
a\alpha\alpha+2b\alpha\gamma-a'\gamma\gamma&=&A&\qquad\mathinner{\ldotp\ldotp\ldotp\ldotp\ldotp\ldotp\ldotp\ldotp}\qquad[1]\\[4pt]
a\beta\beta+2b\beta\delta-a'\delta\delta&=&-A'&\qquad\mathinner{\ldotp\ldotp\ldotp\ldotp\ldotp\ldotp\ldotp}\qquad[2]
\end{array}
\]
From these we obtain
\[
\begin{array}{rcll}
\dfrac{\alpha}{\gamma}&=&\dfrac{\pm\surd\bigl(D+\frac{aA}{\gamma\gamma}\bigr)-b}{{a}}&\qquad\mathinner{\ldotp\ldotp\ldotp\ldotp\ldotp\ldotp\ldotp\ldotp}\qquad[3]\\[12pt]
\dfrac{\beta}{\delta}&=&\dfrac{\pm\surd\bigl(D-\frac{aA'}{\delta\delta}\bigr)-b}{a}&\qquad\mathinner{\ldotp\ldotp\ldotp\ldotp\ldotp\ldotp\ldotp\ldotp}\qquad[4]
\end{array}
\]
\SourceLeaf{182}
\begin{equation*}
\frac{\gamma}{\alpha}=\frac{\pm\surd\bigl(D-\frac{a'A}{\alpha\alpha}\bigr)+b}{a'}
\quad \mathinner{\ldotp\ldotp\ldotp\ldotp\ldotp\ldotp\ldotp\ldotp} \tag*{[5]}
\end{equation*}
\begin{equation*}
\frac{\delta}{\beta}=\frac{\pm\surd\bigl(D+\frac{a'A'}{\beta\beta}\bigr)+b}{a'}
\quad \mathinner{\ldotp\ldotp\ldotp\ldotp\ldotp\ldotp\ldotp\ldotp} \tag*{[6]}
\end{equation*}
Equation 3, 4. 5, 6 must be discarded if $\gamma,\delta,\alpha,\beta$ respectively $=0$. --- But it remains uncertain which \textit{signs} to give the radical quantities. We decide this as follows.

In [3] and [4] the upper signs must clearly be taken when neither $\frac{\alpha}{\gamma}$ nor $\frac{\beta}{\delta}$ has the sign opposite to that of $a$, since taking the lower sign would make $\frac{a\alpha}{\gamma}$ and $\frac{a\beta}{\delta}$ negative quantities. But $A$ and $A'$ have the same sign, so $\surd D$ falls between $\surd\bigl(D+\frac{aA}{\gamma\gamma}\bigr)$ and $\surd\bigl(D-\frac{aA'}{\delta\delta}\bigr)$, and therefore in this case $\frac{\surd D-b}{a}$ lies between $\frac{\alpha}{\gamma}$ and $\frac{\beta}{\delta}$. So the first part of the theorem is proved for the first case.

In the same way we see that in [5] and [6] the lower signs must be taken when neither $\frac{\gamma}{\alpha}$ nor $\frac{\delta}{\beta}$ has the same sign as $a'$ or $a$, since taking the upper sign would necessarily make $\frac{a'\gamma}{\alpha}$, $\frac{a'\delta}{\beta}$ positive quantities. It follows at once that in this case $\frac{-\surd D+b}{a'}$ lies between $\frac{\gamma}{\alpha}$ and $\frac{\delta}{\beta}$. So the second part of the theorem is also proved for the second case. If it could be shown just as easily that in [3] and [4] the lower signs must be taken when neither of the quantities $\frac{\alpha}{\gamma}$, $\frac{\beta}{\delta}$ has the same sign as $a$, and in [5] and [6] the upper signs when neither $\frac{\gamma}{\alpha}$ nor $\frac{\delta}{\beta}$ has the opposite sign, it would follow in the same way that in the former case $\frac{-\surd D-b}{a}$ lies between $\frac{\alpha}{\gamma}$ and $\frac{\beta}{\delta}$, and in the latter $\frac{\surd D+b}{a'}$ between $\frac{\gamma}{\alpha}$ and $\frac{\delta}{\beta}$. The first part of the theorem would then also be proved for the second case, and the second part for the first. But although this is not difficult, it cannot be done without some detours. We prefer the following method.

When none of the numbers $\alpha,\beta.\gamma,\delta=0$, $\frac{\alpha}{\gamma}$ and $\frac{\beta}{\delta}$ have the same signs as $\frac{\gamma}{\alpha}$, $\frac{\delta}{\beta}$. So when neither of these quantities has the same sign as $a'$ or $a$, and therefore $\frac{-\surd D+b}{a'}$ falls between $\frac{\gamma}{\alpha}$ and $\frac{\delta}{\beta}$, neither of the quantities $\frac{\alpha}{\gamma}$, $\frac{\beta}{\delta}$ has the same sign as $a$, and $\frac{a'}{-\surd D+b}=\frac{-\surd D-b}{a}$ (because $aa'=D-bb$) falls between $\frac{\alpha}{\gamma}$ and $\frac{\beta}{\delta}$. Thus, provided neither $\alpha$ nor $\beta=0$, the first part of the theorem is also proved for the second case (for the condition that neither $\gamma$ nor $\delta=0$ is already included in the theorem itself). Similarly, when none of the numbers $\alpha,\beta,\gamma,\delta=0$,
\SourceLeaf{183}
and neither $\frac{\alpha}{\gamma}$ nor $\frac{\beta}{\delta}$ has the sign opposite to that of $a$ or $a'$, so that ${\frac{\surd D-b}{a'}}$ lies between $\frac{\alpha}{\gamma}$ and $\frac{\beta}{\delta}$, then $\frac{\gamma}{\alpha}$ and $\frac{\delta}{\beta}$ also do not have the sign opposite to that of $a'$, and $\frac{a}{\surd D-b}=\frac{\surd D+b}{a'}$ falls between $\frac{\gamma}{\alpha}$ and $\frac{\delta}{\beta}$. Thus, provided neither $\gamma$ nor $\delta=0$, the second part of the theorem is also proved for the second case.

It would therefore remain only to prove that the first part of the theorem also holds in the second case if either of the numbers $\alpha,\beta$ is $=0$, and the second part in the first case if either $\gamma$ or $\delta=0$. But \textit{all these cases are impossible}. For the \textit{first} part of the theorem, suppose that neither $\gamma$ nor $\delta=0$, that $\frac{\alpha}{\gamma}$, $\frac{\beta}{\delta}$ do not have the same sign as $a$, and that 1) $\alpha=0$. Then the equation $\alpha\delta-\beta\gamma=\pm1$ gives $\beta=\pm1$, $\gamma=\pm1$. From [1] we get $A=-a'$, so $A$ and $a'$, and therefore also $a$ and $A'$, have opposite signs, giving $\surd\bigl(D-\frac{aA'}{\delta\delta}\bigr)>\surd D>b$. In [4] the lower sign must therefore be taken, since with the upper sign $\frac{\beta}{\delta}$ would clearly acquire the same sign as $a$. So $\frac{\beta}{\delta}>\frac{-\surd D-b}{a}>1$ (because $a<\surd D+b$ by the definition of a reduced form), \textit{which is absurd}, since $\beta=\pm1$ and $\delta$ is not $=0$. --- 2) Let $\beta=0$. Then the equation $\alpha\delta-\beta\gamma=\pm1$ gives $\alpha=\pm1$, $\delta=\pm1$. From [2] we get $-A'=-a'$, so $a'$ and $a$ and $A$ have the same sign, giving ${\surd\bigl(D+\frac{aA}{\alpha\alpha}\bigr)}>\surd D>b$. In [3] the lower sign must therefore be taken, since with the upper sign $\frac{\alpha}{\gamma}$ would acquire the same sign as $a$. So $\frac{\alpha}{\gamma}>\frac{-\surd D-b}{a}>1$, \textit{which is absurd} for the same reason as before. --- For the \textit{second} part, suppose that neither $\alpha$ nor $\beta=0$, that $\frac{\gamma}{\alpha}$, $\frac{\delta}{\beta}$ do not have the sign opposite to that of $a'$, and that 1) $\gamma=0$. The equation $\alpha\delta-\beta\gamma=\pm1$ gives $\alpha=\pm1$, $\delta=\pm1$. From [1] we get $A=a$, so $a'$ and $A'$ have the same sign, giving $\surd\bigl(D+\frac{a'A'}{\beta\beta}\bigr)>\surd D>b$. In [6] the upper sign must therefore be taken, since with the lower sign $\frac{\delta}{\beta}$ would acquire the sign opposite to that of $a'$. Thus $\frac{\delta}{\beta}>\frac{\surd D+b}{a'}>1$, \textit{which is absurd}, since $\delta=\pm1$ and $\beta$ is not $=0$. Finally, 2) if $\delta=0$, then $\alpha\delta-\beta\gamma=\pm1$ gives $\beta=\pm1$, $\gamma=\pm1$, and therefore from [2], $-A'=a$. This gives $\surd\bigl(D-\frac{a'A}{\alpha\alpha}\bigr)>\surd D>b$, so in [5] the upper sign must be taken. Thus $\frac{\gamma}{\alpha}>\frac{\surd D+b}{a'}>1$, \textit{which is absurd}. --- So the theorem is proved in its full scope.

\Needspace{9\baselineskip}
Since the difference between $\frac{\alpha}{\gamma}$ and $\frac{\beta}{\delta}$ is $=\frac{1}{\gamma\delta}$, the difference between $\frac{\pm\surd D-b}{a}$ and $\frac{\alpha}{\gamma}$ or $\frac{\beta}{\delta}$ is $<\frac{1}{\gamma\delta}$; and between $\frac{\pm\surd D-b}{a}$ and $\frac{\alpha}{\gamma}$, or between that quantity and $\frac{\beta}{\delta}$, no fraction can lie whose denominator is not greater than $\gamma$ or $\delta$ (\textit{the preceding lemma}). --- In the same way the difference of the quantity ${\frac{\pm\surd D+b}{a}}$ from the fraction $\frac{\gamma}{\alpha}$ or $\frac{\delta}{\beta}$
\SourceLeaf{184}
is less than $\frac{1}{\alpha\beta}$, and between that quantity and either of these fractions no fraction can lie whose denominator is not greater than $\alpha$ and $\beta$.

\Needspace{8\baselineskip}
\Article{192}
Applying the preceding theorem to the algorithm of Article~188 shows that the quantity $\frac{\surd D-b}{a}$, which we will denote by $L$, lies between $\frac{\alpha'}{\gamma'}$ and $\frac{\beta'}{\delta'}$, between $\frac{\alpha''}{\gamma''}$ and $\frac{\beta''}{\delta''}$, between $\frac{\alpha'''}{\gamma'''}$ and $\frac{\beta'''}{\delta'''}$, etc. (for it follows easily from the end of Article~189,~3 that none of these bounds has the sign opposite to that of $a$, so the radical quantity $\surd D$ must be given the positive sign), or between $\frac{\alpha'}{\gamma'}$ and $\frac{\alpha''}{\gamma''}$, between $\frac{\alpha''}{\gamma''}$ and $\frac{\alpha'''}{\gamma'''}$, etc. Thus all the fractions $\frac{\alpha'}{\gamma'}$, $\frac{\alpha'''}{\gamma'''}$, $\frac{\alpha'''''}{\gamma'''''}$ etc. lie on the same side of $L$, and all $\frac{\alpha''}{\gamma''}$, $\frac{\alpha''''}{\gamma''''}$, $\frac{\alpha''''''}{\gamma''''''}$ etc. on the other side. But since $\gamma'<\gamma'''$, $\frac{\alpha'}{\gamma'}$ lies \textit{outside} the interval between $\frac{\alpha'''}{\gamma'''}$ and $L$, and for a similar reason $\frac{\alpha''}{\gamma''}$ outside $L$ and $\frac{\alpha''''}{\gamma''''}$, $\frac{\alpha'''}{\gamma'''}$ outside $L$ and $\frac{\alpha'''''}{\gamma'''''}$, etc. So these quantities clearly lie in the following order:
\[
\frac{\alpha'}{\gamma'},\ \frac{\alpha'''}{\gamma'''},\ \frac{\alpha'''''}{\gamma'''''}\mathinner{\ldotp\ldotp\ldotp} L\mathinner{\ldotp\ldotp\ldotp}\frac{\alpha''''''}{\gamma''''''},\ \frac{\alpha''''}{\gamma''''},\ \frac{\alpha''}{\gamma''}
\]
The difference between $\frac{\alpha'}{\gamma'}$ and $L$ is less than the difference between $\frac{\alpha'}{\gamma'}$ and $\frac{\alpha''}{\gamma''}$, \textit{that is}, $<\frac{1}{\gamma'\gamma''}$, and similarly the difference between $\frac{\alpha''}{\gamma''}$ and $L$ is $<\frac{1}{\gamma''\gamma'''}$, etc. So the fractions $\frac{\alpha'}{\gamma'}$, $\frac{\alpha''}{\gamma''}$, $\frac{\alpha'''}{\gamma'''}$ etc. keep approaching the limit $L$ more closely, and since $\gamma',\gamma'',\gamma'''$ keep growing without bound, the difference between the fractions and the limit can be made smaller than any given quantity.

By Article~189, none of the quantities $\frac{\gamma}{\alpha}$, $\frac{{}'\gamma}{{}'\alpha}$, $\frac{{}''\gamma}{{}''\alpha}$ etc. has the same sign as $a$. An argument entirely similar to the preceding one therefore shows that these quantities and $\frac{-\surd D+b}{a'}$, which we will denote by $L'$, lie in the following order:
\[
\frac{\gamma}{\alpha},\ \frac{{}''\gamma}{{}''\alpha},\ \frac{{}''''\gamma}{{}''''\alpha}\mathinner{\ldotp\ldotp\ldotp\ldotp\ldotp} L'\mathinner{\ldotp\ldotp\ldotp\ldotp\ldotp}\frac{{}'''''\gamma}{{}'''''\alpha},\ \frac{{}'''\gamma}{{}'''\alpha},\ \frac{{}'\gamma}{{}'\alpha}
\]
The difference between $\frac{\gamma}{\alpha}$ and $L'$ is less than $\frac{1}{{}'\alpha\alpha}$, the difference between $\frac{{}'\gamma}{{}'\alpha}$ and $L'$ less than $\frac{1}{{}''\alpha{}'\alpha}$, etc. So the fractions $\frac{\gamma}{\alpha}$, $\frac{{}'\gamma}{{}'\alpha}$ etc. keep approaching $L'$ more closely, and the difference can be made smaller than any given quantity.

In the \textit{example} of Article~188 we have $L=\frac{\surd79-{8}}{3}=0{,}2960648$, and the approximating fractions are $\frac{0}{1}$, $\frac{1}{3}$, $\frac{2}{7}$, $\frac{3}{10}$, $\frac{5}{17}$, $\frac{8}{27}$, $\frac{45}{152}$, $\frac{143}{483}$, etc. Now $\frac{143}{483}=0{,}2960662$. --- There we also have $L'=\frac{-\surd79+{8}}{5}=-0{,}1776388$, and the approximating fractions are $\frac{0}{1}$, $-\frac{1}{5}$, $-\frac{1}{6}$, $-\frac{2}{11}$, $-\frac{3}{17}$, $-\frac{8}{45}$, $-\frac{27}{152}$, $-\frac{143}{805}$, etc. But $\frac{143}{805}=0{,}1776397$.
